\clearpage
\section*{Phụ lục: Bảng tra cứu công thức nhanh}
\addcontentsline{toc}{section}{Phụ lục: Bảng tra cứu công thức nhanh}

\subsection*{A. Diện tích hình phẳng}

\begin{hopdinhly}[Bảng công thức]
\renewcommand{\arraystretch}{2.0}
\begin{tabular}{>{\raggedright\arraybackslash}p{5.5cm}>{\raggedright\arraybackslash}p{7.5cm}}
\textbf{Trường hợp} & \textbf{Công thức} \\
\hline
1 đường cong $y=f(x)$ và trục $Ox$, trên $[a;b]$ &
$\displaystyle S=\int_a^b \ababs{f(x)}\dx$ \\
2 đường cong $y=f(x)$, $y=g(x)$, trên $[a;b]$ &
$\displaystyle S=\int_a^b \ababs{f(x)-g(x)}\dx$ \\
Thiếu cận (đường cong khép kín) &
Giải $f(x)=g(x)$ để tìm cận; nếu có $>2$ nghiệm trong miền xét, phải chẻ tích phân theo từng đoạn \\
Parabol -- cát tuyến (công thức nhanh) &
Nếu $f(x)-g(x)=a(x-x_1)(x-x_2)$ thì $\displaystyle S=\frac{\ababs{a}}{6}(x_2-x_1)^3$ \\
\end{tabular}
\end{hopdinhly}

\subsection*{B. Thể tích vật thể và khối tròn xoay}

\begin{hopdinhly}[Bảng công thức]
\renewcommand{\arraystretch}{2.0}
\begin{tabular}{>{\raggedright\arraybackslash}p{5.5cm}>{\raggedright\arraybackslash}p{7.5cm}}
\textbf{Trường hợp} & \textbf{Công thức} \\
\hline
Vật thể biết diện tích thiết diện $S(x)$, trên $[a;b]$ &
$\displaystyle V=\int_a^b S(x)\dx$ \quad (KHÔNG có $\pi$, KHÔNG bình phương) \\
Khối tròn xoay quanh $Ox$, giới hạn bởi $y=f(x)$, $Ox$, trên $[a;b]$ &
$\displaystyle V=\pi\int_a^b f^2(x)\dx$ \\
Khối tròn xoay quanh $Ox$, giữa 2 đường cong $f(x)\ge g(x)\ge0$ &
$\displaystyle V=\pi\int_a^b\ababs{f^2(x)-g^2(x)}\dx$ \\
Khối tròn xoay quanh $Oy$ (đổi vai trò $x\leftrightarrow y$) &
$\displaystyle V=\pi\int_c^d x^2(y)\dif{y}$ \\
\end{tabular}
\end{hopdinhly}

\subsection*{C. Diện tích thiết diện thường gặp (theo cạnh/bán kính $a(x)$)}

\begin{hopdinhly}[Bảng công thức]
\renewcommand{\arraystretch}{1.8}
\begin{tabular}{>{\raggedright\arraybackslash}p{5.5cm}>{\raggedright\arraybackslash}p{7.5cm}}
\textbf{Hình dạng thiết diện} & \textbf{Diện tích $S(x)$} \\
\hline
Hình vuông, cạnh $a(x)$ & $S(x)=a^2(x)$ \\
Tam giác đều, cạnh $a(x)$ & $S(x)=\dfrac{\sqrt3}{4}a^2(x)$ \\
Nửa hình tròn, bán kính $r(x)$ & $S(x)=\dfrac{\pi}{2}r^2(x)$ \\
Hình chữ nhật, hai cạnh $a(x)$, $b(x)$ & $S(x)=a(x)\cdot b(x)$ \\
\end{tabular}
\end{hopdinhly}

\subsection*{D. Công thức hình học không gian đối chiếu (kiểm chứng bằng tích phân)}

\begin{hopdinhly}[Bảng công thức]
\renewcommand{\arraystretch}{1.8}
\begin{tabular}{>{\raggedright\arraybackslash}p{5.5cm}>{\raggedright\arraybackslash}p{7.5cm}}
\textbf{Hình khối} & \textbf{Thể tích} \\
\hline
Khối nón (bán kính đáy $R$, cao $h$) & $V=\dfrac13\pi R^2h$ \\
Khối trụ (bán kính đáy $R$, cao $h$) & $V=\pi R^2h$ \\
Khối cầu (bán kính $R$) & $V=\dfrac43\pi R^3$ \\
\end{tabular}
\end{hopdinhly}

\subsection*{E. Bảng quy đổi đơn vị thường dùng}

\begin{hopdinhly}[Bảng quy đổi]
\renewcommand{\arraystretch}{1.6}
\begin{tabular}{>{\raggedright\arraybackslash}p{7cm}>{\raggedright\arraybackslash}p{6cm}}
\textbf{Quy đổi} & \textbf{Giá trị} \\
\hline
$1\text{ dm}^3$ & $=1$ lít \\
$1\text{ m}^3$ & $=1000$ lít \\
$1\text{ cm}^3$ & $=1$ ml \\
$1\text{ m}^2$ & $=10\,000\text{ cm}^2$ \\
\end{tabular}
\end{hopdinhly}

\subsection*{F. Tổng hợp các sai lầm cần tránh}

\begin{hopchuy}[Tổng hợp sai lầm thường gặp]
\begin{enumerate}[leftmargin=1cm]
  \item \textbf{Quên xét dấu:} Viết thẳng $S=\displaystyle\int_a^b f(x)\dx$ mà không kiểm tra nghiệm của $f(x)=0$ có thuộc $(a;b)$ hay không (Bài 1, Dạng 1).
  \item \textbf{Nhầm diện tích với tích phân có dấu:} $\displaystyle\int_a^b f(x)\dx$ có thể âm hoặc bằng hiệu hai miền, trong khi diện tích luôn là \emph{tổng} các miền dương (Bài 1, Dạng 3).
  \item \textbf{Quên chẻ tích phân} khi hai đường cong cắt nhau nhiều hơn 2 điểm trong miền đang xét (Bài 1, Dạng 2).
  \item \textbf{Nhầm công thức thể tích vật thể với khối tròn xoay:} $V=\displaystyle\int_a^b S(x)\dx$ \textbf{không có $\pi$}, còn $V=\pi\displaystyle\int_a^b f^2(x)\dx$ \textbf{có $\pi$} (Bài 2, Nền tảng).
  \item \textbf{Nhầm $\pi\displaystyle\int\ababs{f^2-g^2}\dx$ với $\pi\displaystyle\int(f-g)^2\dx$} khi tính thể tích khối tròn xoay giữa hai đường cong -- đây là hai đại lượng hoàn toàn khác nhau (Bài 2, Dạng 3).
  \item \textbf{Sai đơn vị} khi chuyển đổi giữa m$^3$, dm$^3$, cm$^3$ và lít, ml trong các bài toán mô hình hóa thực tiễn.
\end{enumerate}
\end{hopchuy}

\newpage
